% --------------------------------------------------------------
% Corporate Twin — Lab Report
% Pre-change Network Digital Twin & Attack-Path Risk Analysis
% --------------------------------------------------------------

\documentclass[11pt,a4paper]{article}
\usepackage[margin=1in]{geometry}
\usepackage[hidelinks]{hyperref}
\usepackage{booktabs}
\usepackage{tabularx}
\usepackage{longtable}
\usepackage{color}
\usepackage{graphicx}
\usepackage{float}
\usepackage{amsmath,amssymb}
\usepackage{enumitem}
\usepackage{array}
\usepackage{calc}

\title{Corporate Twin — Pre-Change Network Digital Twin \& Attack-Path Risk Analysis\\
\vspace{0.5cm} \large Lab Report}
\author{}
\date{\today}

\begin{document}

\maketitle

\begin{abstract}
Corporate Twin is a platform for validating firewall and network security changes
before they reach production. It creates a ``digital twin'' of a branch office
network, captures a baseline of routing, reachability, and firewall policy, and
then simulates candidate changes on a copy of the model. The differential engine
reports exactly what paths open, what close, and how far an attacker could travel,
including a 0--100 risk score, blast-radius footprint, and seven explainable
compliance controls. The platform is designed to run on a Raspberry Pi~5 (4~GB
ARM64) with no Docker dependency when operating in model-only mode, while also
supporting full Docker/Containerlab deployment for live capture.
\end{abstract}

\section{Overview}
Corporate Twin answers the question: ``If I make this change, what actually happens?''
rather than the binary ``is it blocked?'' A change may leave every path reachable while
making one of them vastly more exploitable -- the class of regression this platform is
built to catch.

The core architecture comprises:
\begin{itemize}
  \item \textbf{Snapshot engine} -- captures the baseline (routing, reachability, firewall
        policy) from live containers or from the declarative model.
  \item \textbf{Differential engine} -- applies a candidate change to a \emph{copy} of the
        baseline (never mutating the original) and computes policy deltas, reachability
        deltas, and risk deltas.
  \item \textbf{Risk engine} -- blends four explainable components (blast radius 40\%,
        zone penetration 25\%, crown-jewel exposure 25\%, path novelty 10\%) into a
        composite score with buckets: HIGH $\ge$ 70, MEDIUM $\ge$ 30, LOW $<$ 30.
  \item \textbf{Compliance controls} -- seven named rules (CTL-001 through CTL-007)
        evaluating default-deny posture, unrestricted permits, crown-jewel cross-zone
        exposure, management-plane reachability, blanket permits into crown-jewel zones,
        and legacy cleartext protocols.
  \item \textbf{Safe rollback} -- the baseline is immutable; rollback is logical and
        trustworthy.
\end{itemize}

\section{Project Status}
The project is \textbf{complete}. Evidence:

\begin{tabularx}{\linewidth}{@{} l X @{}}
  \toprule
  \textbf{Aspect} & \textbf{Status} \\
  \midrule
  Tests & 101 tests passed (0 skipped) \\
  Linting & \texttt{ruff check} all checks passed \\
  Type checking & \texttt{mypy} success on 35 files \\
  API & Full FastAPI backend with \texttt{/api/v1} endpoints \\
  Documentation & Complete README with Quickstart, API reference, design decisions \\
  Modes & Both Docker/Containerlab live mode and model-only mode \\
  Error handling & Standardized error envelopes with machine-readable codes \\
  Configuration & Environment variables with repository-relative defaults \\
  \bottomrule
\end{tabularx}

All 101 tests pass without a Docker daemon; the runtime is faked at the process
boundary, exercising real node validation, container-state checks, capture logic, and
command timeouts in model-only mode.

\section{Features}
\begin{description}
  \item [Change types accepted by \texttt{simulate-change}]:
    \texttt{allow\_flow}, \texttt{deny\_flow}, \texttt{add\_policy\_rule},
    \texttt{remove\_policy\_rule}, \texttt{set\_link\_state}, \texttt{compromise\_node}.
  \item [Risk scoring] blends four components:

    \noindent\begin{tabularx}{\linewidth}{>{\hsize=0.8\hsize}X >{\hsize=0.4\hsize}X >{\hsize=1.8\hsize}X}
      \toprule
      \textbf{Component} & \textbf{Weight} & \textbf{Meaning} \\
      \midrule
      Blast radius & 0.40 & Weighted share of the estate reachable from the origin \\
      Zone penetration & 0.25 & Fraction of zones beyond the origin's own zone that were reached \\
      Crown-jewel exposure & 0.25 & Exploitability of the most exposed crown-jewel asset \\
      Path novelty & 0.10 & Share of previously blocked origin paths this change opens \\
      \bottomrule
    \end{tabularx}

    Buckets: HIGH $\ge$ 70, MEDIUM $\ge$ 30, LOW $<$ 30. Every weight and threshold
    is a named constant in \texttt{backend/app/core/risk.py}.
  \item [Exposure widening] A port-agnostic path that widens to every port is flagged
        even though it stayed reachable (the \texttt{exposure\_widened} signal).
  \item [Compliance] Seven controls are evaluated on every baseline and every candidate
        change, each carrying severity, observed evidence, and remediation text.
  \item [Modes] Runs fully without Docker/Containerlab (model-only) or with live lab
        deployment. Health endpoint reports \texttt{status: degraded} when the lab is
        unavailable.
\end{description}

\section{Testing}
The test suite (101 tests) covers:
\begin{itemize}
  \item \texttt{iptables} parsing, first-match-wins evaluation, port-interval algebra;
  \item multi-hop routing, link-state changes, port-set intersection;
  \item risk scoring, blast radius and the differential engine;
  \item full end-to-end HTTP tests through the ASGI application, including error
        envelopes and baseline persistence across a simulated restart.
\end{itemize}

Verification command:
\begin{verbatim}
.venv/bin/ruff check backend frontend scripts && \
.venv/bin/mypy && \
.venv/bin/python -m pytest
\end{verbatim}

\section{Quickstart}
\begin{verbatim}
# Set up the environment
git clone <this-repo> && cd corporate-twin
scripts/setup_env.sh        # creates .venv & installs dependencies

# Run the backend (model-only mode; no Docker required)
scripts/run_backend.sh        # -> http://127.0.0.1:8000

# Run the dashboard
scripts/run_frontend.sh       # -> http://127.0.0.1:8501

# Capture baseline and validate a change
curl -s -X POST http://localhost:8000/api/v1/snapshot | \
  jq '.snapshot | {snapshot_id, runtime_mode, reachable_count}'

curl -s -X POST http://localhost:8000/api/v1/simulate-change \
  -H 'Content-Type: application/json' \
  -d '{"change_type":"allow_flow","target_node":
       "perimeter-gateway","chain":"FORWARD",
       "src":"10.10.10.0/24","dst":"10.10.20.50",
       "protocol":"tcp","description":"CHG-1042"}' | \
  jq '.risk.projected | {risk_score, risk_level, blast_radius_pct}'
\end{verbatim}

\section{Raspberry Pi Compatibility}
\textbf{Yes -- the project is designed for a Raspberry Pi~5 with 4~GB RAM.}

\begin{tabularx}{\linewidth}{>{\hsize=0.6\hsize}X >{\hsize=1.4\hsize}X}
  \toprule
  \textbf{Feature} & \textbf{Pi~5 / 4~GB relevance} \\
  \midrule
  Architecture & arm64-optimized; containers are Alpine + \texttt{iproute2} + \texttt{iptables} ($\sim$25~MB each) \\
  Python & Targets Python 3.11 deliberately -- ``Raspberry Pi OS Bookworm ships Python 3.11'' \\
  Mode & Runs \textbf{fully} in model-only mode with no Docker daemon required. The backend validates changes against \texttt{config/model.yaml} only. \\
  Resource footprint & ``No hypervisors, no VM-backed network OS images, no GNS3/EVE-NG dependency. The backend keeps working with the lab torn down, so a laptop with no Docker daemon can still validate changes against the model.'' (README) \\
  Lab deployment (optional) & If live containers are desired, \texttt{scripts/build\_images.sh} and \texttt{scripts/deploy.sh} work on Pi~5 -- they just take longer. \\
  \bottomrule
\end{tabularx}

\textbf{Bottom line:} The core solver, risk engine, diff engine, and all API endpoints work
without Docker by reading from \texttt{config/model.yaml}. Install the environment and
try \texttt{scripts/run\_backend.sh} -- it will start the FastAPI server in model-only mode
immediately on the Pi with no Docker needed.

\section{Initial Commit Message}
Since this repository has no commits yet, a natural initial commit message is:

\begin{verbatim}
Initial commit — Corporate Twin: pre-change network
digital twin & attack-path risk analysis

A platform for validating firewall/security changes before
production. Captures network baseline, simulates candidate
changes, and reports what opens/closes/widens, with a
0--100 risk score and 7 explainable compliance controls.
\end{verbatim}

\section{Gotchas \& Scope Limits}
\begin{itemize}
  \item \texttt{model.yaml} $\leftrightarrow$ \texttt{corporate-twin.yml} must agree (node names, interface
        names, IPs). The loader refuses to start on mismatch.
  \item No importer for real device configs; adapting to a real network = hand-transcribing
        into \texttt{model.yaml}.
  \item Rules are first-match-wins; the conntrack \texttt{ESTABLISHED,RELATED} rule must stay
        first in \texttt{FORWARD}.
  \item Crown jewel \texttt{database-server} is deliberately default-ALLOW on INPUT -- models a
        common real-world latent risk (trips CTL-002).
  \item IPv4 forwarding, static routing, \texttt{INPUT}/\texttt{FORWARD}/\texttt{OUTPUT} only.
        No NAT, dynamic routing (OSPF/BGP), user-defined chains, VLANs, or IPv6.
  \item The blast-radius model estimates reachable \emph{exposure}; it does not prove
        exploitability.
\end{itemize}

\clearpage
\appendix
\section{Reference API Endpoints (all under \texttt{/api/v1})}
\begin{tabularx}{\linewidth}{@{} l l X @{}}
  \toprule
  \textbf{Method} & \textbf{Path} & \textbf{Purpose} \\
  \midrule
  POST & \texttt{/snapshot} & Capture the pre-change baseline \\
  POST & \texttt{/simulate-change} & Apply a candidate change or simulate a compromised node \\
  POST & \texttt{/analyze-blast-radius} & Graph-based reachability analysis with risk score and blast radius \\
  POST & \texttt{/rollback} & Revert the network state to the pre-change baseline \\
  GET & \texttt{/health} & Readiness, runtime mode and named checks \\
  GET & \texttt{/topology} & Nodes, links and zones for map rendering \\
  GET & \texttt{/policies} & Compiled chains with defaults and rules in order \\
  GET & \texttt{/compliance} & Compliance score and per-control findings \\
  \bottomrule
\end{tabularx}

\end{document}
